\section{Supplementary Methods and Results}

\subsection{Network inventory}

\begin{table}[htbp]
\caption{Source and external target network inventory parsed from the SUMO
\texttt{net.xml} files.  Jinan uses one geometry and three demand files.}
\label{tab:network-scope}
\centering
\scriptsize
\resizebox{\textwidth}{!}{\begin{tabular}{llrrrl}
\toprule
Role & City group & Networks & TLS & Lanes & Scenarios \\
\midrule
Source & atlanta & 1 & 5 & 55 & atlanta\_1x5 \\
Source & cologne & 3 & 12 & 252 & cologne1, cologne3, cologne8 \\
Source & hangzhou & 6 & 36 & 532 & hangzhou\_bc\_tyc, hangzhou\_kn\_hz, hangzhou\_qc\_yn, hangzhou\_sb\_sx, hangzhou\_4x4, hangzhou\_4x4\_hetero \\
Source & ingolstadt & 3 & 32 & 2227 & ingolstadt1, ingolstadt7, ingolstadt21 \\
Source & new\_york & 1 & 196 & 2562 & manhattan\_28x7 \\
Source & resco\_synthetic & 2 & 32 & 360 & grid4x4, arterial4x4 \\
Source & salt\_lake\_city & 2 & 4 & 88 & saltlake\_400s\_200w\_q1\_weekday\_peak, saltlake\_state\_university\_q1\_weekday\_peak \\
External target & los\_angeles & 1 & 4 & 157 & la\_1x4 \\
External target & jinan & 1 & 12 & 186 & jinan\_3x4 (3 demand files) \\
\bottomrule
\end{tabular}
}
\end{table}

The external CityFlow files were converted to SUMO and revised through a
semantic admission audit.  Phase identities use final SUMO connection indices,
permissive links remain distinct from protected movements, all routes are
contiguous and demand is conserved.  A 3,600-s admission run had no teleports
or collisions.  The admitted conversion tree has SHA-256
\texttt{4a2e17618c2ac584fc767562f3a7096c95d8f904f03ef6baa9cd5d34d282eadd}.

\subsection{Sequential development and confirmation}

\begin{table}[htbp]
\caption{Complete high-level decision trail for the pairwise method.  Positive
percentages denote lower normalized action regret than the rigid anchor.  Seed
7079 is diagnostic evidence that informed v43 and is not reused as v43
confirmation.}
\label{tab:development-history}
\centering
\scriptsize
\resizebox{\textwidth}{!}{\begin{tabular}{llll}
\toprule
Stage & Independent check & Result & Protocol consequence \\
\midrule
v38 global pairwise & Salt Lake, 2 networks & $-12.38\%$ vs rigid; FAIL & reject global correction \\
v39 anchored selector & 7-city development & $-0.95\%$ macro; FAIL & reject first anchor blend \\
v40 cross-fitted selector & 7-city nested folds & selected anchor; FAIL & add seed-held-out evaluation \\
v41 seed-held-out & seed 4047, 7 cities & $+9.84\%$, 6/7 improve; PASS & freeze for external diagnosis \\
v42 external diagnostic & seed 7079, 2 cities & $+4.50\%$ macro, 1/2 improve; FAIL & exclude seed; remove B100 cap \\
v43 external confirmation & seed 8171, 2 cities & $+8.09\%$, 2/2 improve; PASS & authorize closed loop \\
\bottomrule
\end{tabular}
}
\end{table}

The v42 failure is retained because its city-macro improvement alone would have
looked positive.  It failed the predeclared requirement that both external
cities improve: Los Angeles regressed by 0.00072 normalized regret.  V43 then
removed the arbitrary 100-group cap, used every valid adaptation group and
cross-fitted by complete adaptation seed.  The later seed 8171, rather than
7079, is the formal external confirmation.

\subsection{City-resolved closed-loop outcomes}

\begin{table}[htbp]
\caption{External closed-loop outcomes before equal-city aggregation.  Waiting
is the all-departed 3,600-s metric in seconds.  Incident entries are means over
the two closed-loop seeds within that city and policy.}
\label{tab:external-city-results}
\centering
\scriptsize
\resizebox{\textwidth}{!}{\begin{tabular}{lrrrrrr}
\toprule
Policy & LA wait & Jinan wait & LA completion & Jinan completion & LA incidents & Jinan incidents \\
\midrule
Fixed time & 507.33 & 169.84 & 0.449 & 0.873 & 4.00 & 0.00 \\
Dense residual & 212.88 & 105.41 & 0.838 & 0.902 & 0.50 & 0.00 \\
Simulator only & 219.17 & 94.26 & 0.851 & 0.903 & 0.00 & 0.00 \\
Rigid anchor & 206.78 & 99.05 & 0.862 & 0.902 & 1.00 & 0.00 \\
\textbf{Anchored CFCMT} & 204.42 & 93.91 & 0.874 & 0.903 & 0.00 & 0.00 \\
Phase pressure & 197.51 & 89.73 & 0.890 & 0.905 & 0.50 & 0.00 \\
MaxPressure & 191.89 & 91.58 & 0.888 & 0.904 & 0.50 & 0.00 \\
\bottomrule
\end{tabular}
}
\end{table}

The CFCMT-versus-dense waiting improvements are 8.46 s in Los Angeles and
11.50 s in Jinan.  CFCMT-versus-MaxPressure differences are $-12.53$ and
$-2.33$ s, where negative values indicate that MaxPressure is better.  The
opposite directions cannot be hidden by equal-city aggregation.

\subsection{Fresh source-contribution diagnostics}

The comparator artifacts were created in the v43 baseline freeze, before the
original closed-loop outcomes.  Their deployed score assigns zero coefficient
to source-labelled samples and their selected target weights were within
numerical precision of one in both cities.  However, a retained source-trained
prior can resolve exact prediction ties.  We therefore call this the legacy
near-target-only comparator rather than a strict zero-source-row arm.  V90
added these artifacts to the eight original scenario--seed cells only after
the main comparison and is therefore not a confirmation.  V91 froze 64 unique
seeds before execution, required 512 exact policy identities and completed all
of them without seed exclusion.

The largest v91 failure occurred at seed 7417817 in Los Angeles.  CFCMT waiting
was 975.27 s with completion 0.270, compared with 194.74 s and completion 0.919
for the legacy near-target-only comparator; neither policy teleported or
recorded a collision in that cell.  This record is retained in the primary
estimate.  The failure was not solely an outlier artifact: across the 64 Los
Angeles cells, the median relative difference was $+2.75\%$, only 22 favoured
CFCMT and 21 regressed by more than 5\%.  Across 192 Jinan scenario--seed cells,
179 favoured CFCMT, and all 64 Jinan seed-level city contrasts favoured CFCMT.
The city-mean waiting effects were $-3.045\%$ in Jinan and $+7.865\%$ in Los
Angeles, where negative values favour CFCMT.  These diagnostics explain the
opposite city means but do not change the preregistered seed-level analysis.

\subsection{Post-v91 V98 controller-pair confirmation}

V98 is a separate post-v91 protocol and is not pooled with the v91 estimate.
Its target-offline selector returned exact target-only control for Los Angeles
and selected a Hangzhou source component at mass one for Jinan before any fresh
v98 rollout was opened.  Los Angeles was therefore an ex ante abstention, not
a city removed after observing the confirmation outcomes.  The Jinan matrix
used 56 further seeds, three demand files per seed and two controller arms, for
336 rollouts in total.  Relative to the strict zero-source-row target model,
the admitted controller reduced waiting by 4.527\% (paired 95\% interval,
4.279--4.787\%); all 56 seed-level contrasts were favourable
($p=7.5475\times10^{-11}$, Wilcoxon signed-rank test).  The two frozen protocols
are juxtaposed, rather than combined, in
\Cref{tab:jinan-conditional-source-transfer}.

This confirmation supports a protocol-specific Jinan controller-pair gain.  The
preselected controller outperformed its declared strict zero-source-row target
comparator, but that comparator has lower capacity.  V98 therefore does not
isolate source-row contribution at fixed architecture.  Both arms also share the
risk multiplier of 0.5 inherited from earlier Jinan closed-loop development.
The result therefore does not establish a general multicity selector or
superiority over phase pressure.

\subsection{Seven-city source-gate development}

Three later source gates were frozen and evaluated on the seven development
city groups.  The robust-history gate selected a source for six targets, and
the action-agreement gate selected a source for all seven; both improved the
equal-city mean but failed the target-only safety comparison because their
one-sided upper bounds were positive and their largest city regressions
exceeded 0.02 normalized regret.  A stricter mechanism-compatibility gate used
five B20-to-B5 target-label folds to score queue propagation, served movement,
red accumulation, spillback and mobility.  It admitted no source and returned
exact target only in all seven cities (\Cref{tab:source-gate-development}).

\begin{table}[htbp]
\caption{Frozen seven-city source-gate development.  Negative effects favour
the selected gate.  ``Upper'' is the one-sided 95\% upper bound versus the
architecture-matched target-only causal arm.  Dense denotes the pooled
H2O$^+$-style residual.  These are development results, not external
confirmation.}
\label{tab:source-gate-development}
\centering
\scriptsize
\resizebox{\textwidth}{!}{\begin{tabular}{llrrrrrr}
\toprule
Version & Source gate & Source & $\Delta$ target & Upper & Cities & Max. reg. & $\Delta$ dense \\
 &  & selected & mean & 95\% & improved &  & mean \\
\midrule
v145 & Robust history & 6/7 & -0.0028 & +0.0093 & 4/7 & +0.0281 & -0.0720 \\
v146 & OOF action agreement & 7/7 & -0.0016 & +0.0095 & 3/7 & +0.0262 & -0.0708 \\
v148 & Mechanism compatibility & 0/7 & +0.0000 & +0.0000 & 0/7 & +0.0000 & -0.0691 \\
\bottomrule
\end{tabular}
}
\end{table}

The progression exposes a safety--usefulness trade-off rather than a successful
source selector.  v145 and v146 retained source signal but allowed harmful city
tails.  v148 removed those tails only by abstaining everywhere.  Its 7/7
advantage over the dense residual is therefore inherited from the rigid
target-only causal model and cannot be attributed to source-city labels or to a
full mechanism world model.  No external confirmation of these v145--v148
seven-city gate rules was authorized after their failures.

We subsequently tested the alternative explanations raised by this result
without weakening the frozen safety or source-identity criteria
(\Cref{tab:source-remedy-ladder}).  Leave-one-seed diagnostics confirmed that
beneficial source headroom exists after the fact.  Mechanism-block priors,
budget-aware shrinkage, local right-of-way context, state-conditioned utility,
leave-city-out meta-utility and hierarchical priors nevertheless failed to
turn that headroom into a stable deployment rule.  Most conservative variants
returned exact rigid fallback.  The one state-conditioned gate that passed its
offline development screen caused a 37.9\% Cologne regression in the
pre-expansion closed-loop smoke, so its full matrix was not launched.
The later v123/v157a budget contrasts were also retained only as historical
model-path comparisons after a Jinan domain-handling difference was found.
V157B corrected B100 alone and passed its relative source-versus-target gate on
the reused selector, while remaining worse than phase pressure.  The authorized
V157C native one-action diagnostic completed seeds 80314 and 88625, but retained
seed 27178 as invalid because the fixed 480-s checkpoint contained no eligible,
scorable traffic light.  Its two-valid-seed descriptive means were $+0.00236$
for source minus target, $-0.00232$ for source minus matched placebo, $+0.00125$
for target minus PhasePressure and $+0.00362$ for source minus PhasePressure.
Only one of two seeds improved versus target and neither improved versus the
PhasePressure heuristic; all recorded collision counts were zero.  The fixed
three-seed protocol is incomplete and supports no three-seed inference.
V158 then froze the V157B predictors and collected 100 native-prefix groups
(800 candidate rows and 700 non-reference branches) from ten fixed Jinan
development seeds, with no simulator-state restoration.  Five whole-seed folds
gave source-minus-target $-0.00113$ (paired 95\% interval
$[-0.00240,+0.00019]$; 6/10 seeds improved), source-minus-placebo $-0.00072$
($[-0.00219,+0.00091]$; 7/10) and source-minus-PhasePressure $+0.00014$
($[-0.00239,+0.00261]$; 5/10).  Every preregistered OOF comparison failed;
the reserve was therefore neither authorized nor run.  The failure closes this
Jinan native-prefix ranking-remediation route and supports no controller
adoption, unseen-city claim or post-hoc retuning.

\begin{table}[htbp]
\caption{Post-v148 source-remedy ladder.  Effects are normalized action-regret
differences unless seconds or percentages are shown; negative values favour
the source-aware arm.  B100 denotes 100 target adaptation action groups.  The
V157C row reports descriptive means over two valid seeds because the fixed
three-seed protocol was incomplete.  The V158 row reports the preregistered
ten-seed OOF estimates; its gate failed and the reserve was withheld.  The
versioned evidence identifiers are retained in the machine-readable source
data, not in the displayed table.}
\label{tab:source-remedy-ladder}
\centering
\footnotesize
\begin{tabular}{@{}>{\raggedright\arraybackslash}p{0.14\textwidth}>{\raggedright\arraybackslash}p{0.18\textwidth}>{\raggedright\arraybackslash}p{0.09\textwidth}>{\raggedright\arraybackslash}p{0.29\textwidth}>{\raggedright\arraybackslash}p{0.16\textwidth}@{}}
\toprule
Stage & Deployment information & Source & Held-out result & Adjudication \\
\midrule
Repeatability diagnostic & Other development seeds; forced source & Oracle only & 16/21 seed units and 7/7 city means improved & Headroom exists; not deployable \\
Mechanism parameter prior & Target folds plus source mechanism block & 7/7 prototype & 6/7 improved; mean -0.00409; placebo upper 95\% +0.00089 & Offline pass; source identity unresolved \\
Budget-aware shrinkage & B100 target folds; inverse-budget prior & 1/7 & mean -0.00523; no city regression & Reject multicity generality \\
Local mechanism context & B100 plus local right-of-way context & 0/7 & Exact rigid fallback in all cities & Reject \\
State-conditioned utility & B100 nested state-level utility labels & 3/7 & pre-correction historical: offline mean -0.00052; smoke: Cologne +37.9\%, New York +2.6\%, RESCO -17.7\% & Historical; reject after closed-loop smoke \\
Feature-aligned dense utility & V150O replay; stored-name rigid binding & 1/7 inner; 0 follow-up & macro +0.00135 vs rigid; Hangzhou +0.00942/+0.00061/+0.00093 vs rigid/placebo/blind & Reject globally and city-specifically \\
Leave-city-out meta utility & Six pseudo-target cities; target B100 & 0/7 & Exact rigid fallback in all cities & Reject \\
Source-only hierarchical prior & Equal-city source prior; no target gate labels & 0/7 & forced mean +0.00277; 3/7 improved & Reject \\
Target-calibrated hierarchical prior & B100 target selector; same prior grid & 0/7 & forced mean -0.01213, 95\% [-0.03879, +0.00266] & Reject; exact rigid fallback \\
Domain-aligned B100 source gate & V123/V157A confounded; reused 22-seed Jinan selector & One-action branch only & source-target -0.01357, 95\% [-0.01784, -0.00935]; 20/22 improve; +0.02534 vs PP & B100 gate pass; V157C incomplete \\
Native one-action branch & Fixed Jinan roster; seed 27178 t=480 invalid & 2/3 valid seeds & descriptive source-target +0.00236 (1/2 improve); source-placebo -0.00232; target-PP +0.00125; source-PP +0.00362 (0/2 improve); collisions 0; PP=PhasePressure heuristic & Incomplete; no three-seed inference \\
Native-prefix ranking calibration & Frozen V157B predictors; native labels; five seed-held-out folds & OOF 0/3; no reserve & source-target -0.00113, 95\% [-0.00240, +0.00019] (6/10 improve); source-placebo -0.00072, 95\% [-0.00219, +0.00091] (7/10 improve); source-PP +0.00014, 95\% [-0.00239, +0.00261] (5/10 improve); 100 groups/800 rows/700 branches; 0 restores & Reject; OOF gate failed; reserve withheld \\
\bottomrule
\end{tabular}

\end{table}

\subsection{Metric population and safety accounting}

SUMO \texttt{tripinfo} was written with unfinished records at the fixed
3,600-s horizon.  The primary waiting-time population includes all departed
vehicles, with unfinished waiting accrued to the horizon.  Completed-only
waiting and completion ratios are secondary outcomes.  System vehicle-hours
integrate active plus pending vehicles and therefore include demand not yet
inserted.  Across the complete matrix, 171 raw collision events collapsed to
13 unique participant--lane--time incidents in six rollouts; CFCMT had zero.
All starting and ending teleport counters were zero.

\subsection{Discarded pressure-guard development track}

Before the anchored pairwise method was developed, an exploratory controller
placed a learned mechanism residual behind a MaxPressure fallback.  On eight
RESCO scenarios and 600-s rollouts its mean queue was 12.918, close to
MaxPressure at 12.907.  At 3,600 s, phase pressure was materially stronger
(32.234 versus 43.511).  These experiments motivated treating pressure as a
hard performance ceiling, but they use a different controller, target set and
metric from v43 and are excluded from the main confirmation.

\begin{table}[htbp]
\caption{Exploratory 600-s pressure-guard stress test retained for development
transparency.  This is not the final anchored CFCMT controller.}
\label{tab:discarded-guard-600}
\centering
\scriptsize
\resizebox{\textwidth}{!}{\begin{tabular}{lrrrrrrr}
\toprule
Policy & Mean queue & Seed std & P90 queue & Trip time & Trip wait & Time loss & Throughput \\
\midrule
MaxPressure & 12.9068 & 0.1546 & 19.5279 & 87.5891 & 8.0605 & 24.5171 & 0.8267 \\
CFCMT + pressure guard & 12.9180 & 0.1919 & 19.5973 & 87.4863 & 8.0114 & 24.4127 & 0.8256 \\
Phase pressure & 13.2446 & 0.1816 & 19.3911 & 86.5371 & 7.4544 & 23.4628 & 0.8199 \\
H2O+-style dense & 17.2279 & 3.6084 & 28.9325 & 86.6585 & 8.5768 & 24.5488 & 0.7653 \\
Target-static simulator & 20.4023 & 0.1758 & 32.5538 & 97.9633 & 19.9032 & 35.3928 & 0.7429 \\
Native actuated SUMO & 30.1825 & 0.3536 & 47.6873 & 119.7813 & 33.0017 & 57.0805 & 0.7580 \\
Fixed RESCO program & 42.6732 & 0.2709 & 69.6852 & 142.3116 & 54.6333 & 79.8759 & 0.7127 \\
\bottomrule
\end{tabular}
}
\end{table}

\begin{table}[htbp]
\caption{Exploratory 3,600-s pressure-guard boundary check.  Phase pressure is
the strongest policy in this discarded protocol.}
\label{tab:discarded-guard-3600}
\centering
\scriptsize
\resizebox{\textwidth}{!}{\begin{tabular}{lrrrrrrr}
\toprule
Policy & Mean queue & Seed std & P90 queue & Trip time & Trip wait & Time loss & Throughput \\
\midrule
Phase pressure & 32.2343 & 0.9901 & 62.1800 & 104.2291 & 20.8528 & 38.5434 & 0.9492 \\
MaxPressure & 42.7872 & 0.7314 & 71.6277 & 106.9182 & 20.5521 & 40.7331 & 0.9448 \\
CFCMT + pressure guard & 43.5106 & 0.5699 & 74.2496 & 108.6326 & 22.0400 & 42.4485 & 0.9448 \\
Target-static simulator & 44.8973 & 1.2335 & 65.4094 & 258.2547 & 169.8900 & 188.9974 & 0.8754 \\
H2O+-style dense & 49.5058 & 3.9510 & 74.9880 & 221.4173 & 131.1706 & 150.3275 & 0.8827 \\
Native actuated SUMO & 58.3468 & 1.1370 & 99.4055 & 164.5754 & 66.1326 & 98.1604 & 0.9714 \\
Fixed RESCO program & 84.6088 & 0.6458 & 133.0312 & 221.2093 & 113.3988 & 154.3776 & 0.9385 \\
\bottomrule
\end{tabular}
}
\end{table}

\subsection{Post-freeze fresh-seed execution-layer stress test}

The v86--v89 successor track uses a target-state-conditioned 120-s action
originator and global execution cooldown rather than the v43 anchored all-pair
controller.  Negative deltas favour the learned executor relative to phase
pressure.  The first three rows reuse the same 54 disclosed adaptive seeds and
are not independent replications.  Only v89 uses 64 newly generated seeds; its
joint gate was frozen before those rollouts and failed.  The eight v85
prospective seeds were not opened.

\begin{table}[htbp]
\caption{Post-freeze Jinan execution-layer stress test.  Intervals are paired
simulator-seed bootstrap intervals.  Development rows are shown to expose the
adaptive-to-confirmatory shrinkage and are not pooled with v89.}
\label{tab:postfreeze-stress}
\centering
\scriptsize
\resizebox{\textwidth}{!}{\begin{tabular}{llrrrrl}
\toprule
Stage & Evidence & Seeds & Mean $\Delta$ & 95\% CI & Improved & Decision \\
\midrule
v86 legacy global & adaptive screen & 54 & -0.429\% & [-0.818, -0.033]\% & 68.5\% & development only \\
v87 unbounded stay & adaptive ablation & 54 & 0.104\% & [-0.295, 0.519]\% & 46.3\% & variant rejected \\
v88 one-stay budget & adaptive ablation & 54 & -0.242\% & [-0.674, 0.195]\% & 55.6\% & variant rejected \\
v89 frozen legacy & fresh confirmation & 64 & -0.043\% & [-0.449, 0.370]\% & 53.1\% & confirmation rejected \\
\bottomrule
\end{tabular}
}
\end{table}

\subsection{Native-protocol reinforcement-learning checks}

\begin{table}[htbp]
\caption{Official RESCO full-budget RL results under the upstream native
protocol.  Unsupported upstream scenario--algorithm pairs are not silently
patched into new methods.  These rows are not same-protocol comparisons with
CFCMT.}
\label{tab:official-resco-rl}
\centering
\scriptsize
\resizebox{\textwidth}{!}{\begin{tabular}{llrrrrr}
\toprule
Algorithm & Scenario & Runs & Time loss & Duration & Waiting & Elapsed (s) \\
\midrule
IDQN & arterial4x4 & 3 & 1505.536 & 501.305 & 436.920 & 17179.205 \\
MPLight & arterial4x4 & 3 & 1056.797 & 357.454 & 238.745 & 13576.024 \\
IDQN & cologne1 & 3 & 612.737 & 164.905 & 138.950 & 812.441 \\
MPLight & cologne1 & 3 & 189.393 & 122.751 & 68.843 & 6985.784 \\
IDQN & cologne3 & 3 & 86.610 & 92.313 & 43.268 & 2351.500 \\
MPLight & cologne3 & 3 & 1011.239 & 543.511 & 515.329 & 10132.753 \\
IDQN & cologne8 & 3 & 24.760 & 89.379 & 8.354 & 5946.643 \\
IDQN & grid4x4 & 3 & 47.752 & 158.215 & 25.205 & 11120.623 \\
MPLight & grid4x4 & 3 & 44.393 & 156.112 & 16.598 & 10392.985 \\
IDQN & ingolstadt1 & 3 & 18.097 & 37.191 & 5.490 & 911.332 \\
MPLight & ingolstadt1 & 3 & 401.895 & 85.906 & 63.192 & 7888.143 \\
IDQN & ingolstadt21 & 3 & 256.588 & 353.111 & 176.405 & 13141.256 \\
IDQN & ingolstadt7 & 3 & 51.869 & 78.827 & 17.357 & 5313.646 \\
MPLight & ingolstadt7 & 3 & 284.391 & 111.560 & 57.463 & 8755.139 \\
\bottomrule
\end{tabular}
}
\end{table}

\begin{table}[htbp]
\caption{Patched LibSignal full-budget execution audit under native LibSignal
states, rewards, timing and metrics.}
\label{tab:libsignal-rl}
\centering
\scriptsize
\resizebox{\textwidth}{!}{\begin{tabular}{llrrrrrr}
\toprule
Agent & Network & Runs & Travel time & Queue & Delay & Throughput & Elapsed (s) \\
\midrule
dqn & sumo1x1 & 3 & 38.945 & 3.240 & 0.283 & 1998.667 & 175.967 \\
frap & sumo1x1 & 3 & 41.923 & 24.461 & 0.460 & 1664.667 & 439.620 \\
mplight & sumo1x1 & 3 & 40.546 & 4.373 & 0.319 & 1995.000 & 371.569 \\
presslight & sumo1x1 & 3 & 39.748 & 3.442 & 0.285 & 1998.667 & 175.835 \\
dqn & sumo1x3 & 3 & 97.090 & 6.531 & 0.244 & 2487.333 & 425.384 \\
frap & sumo1x3 & 3 & 118.712 & 11.672 & 0.398 & 2505.333 & 1583.053 \\
mplight & sumo1x3 & 3 & 66.502 & 1.876 & 0.243 & 2814.333 & 647.664 \\
presslight & sumo1x3 & 3 & 63.188 & 1.164 & 0.188 & 2816.000 & 393.414 \\
dqn & sumo4x4 & 3 & 140.610 & 0.603 & 0.041 & 1465.000 & 2140.773 \\
frap & sumo4x4 & 3 & 140.477 & 0.612 & 0.042 & 1464.000 & 7188.185 \\
mplight & sumo4x4 & 3 & 140.469 & 0.607 & 0.042 & 1464.333 & 1017.297 \\
presslight & sumo4x4 & 3 & 144.279 & 0.711 & 0.048 & 1463.333 & 2080.043 \\
\bottomrule
\end{tabular}
}
\end{table}

\subsection{Frozen artifact chain}

\begin{table}[htbp]
\caption{SHA-256 chain for the final external result.  The paper artifact
builder verifies all result gates before regenerating this table.}
\label{tab:artifact-chain}
\centering
\tiny
\resizebox{\textwidth}{!}{\begin{tabular}{ll}
\toprule
Artifact & SHA-256 \\
\midrule
development audit & \texttt{\detokenize{5bcae9467cdb4cc157676e94078d204702ebdca6d51007ffca4bc85d13951066}} \\
source cache & \texttt{\detokenize{14c6dfad87999e55fa86fa942ac0ca2334129f4d1662b3c31a6e4fc7aa74c242}} \\
target adaptation cache & \texttt{\detokenize{beb6a968c43dd58ceac4df6d9da8dfca47fc1f0869dfda0dc151a3e9cae58d77}} \\
joint method freeze & \texttt{\detokenize{471e8c34eb7be96411c812c7fe961d5c07f1857e3fb50d2c8d7edd59fc23c596}} \\
held-out seed-8171 result & \texttt{\detokenize{20d682b48bca0f8821ecf98e54b04c94af87d530add51cc1d92d9a176f74cf60}} \\
closed-loop aggregate & \texttt{\detokenize{be36c73c5835e632231140fdbdabc4bba422f54b503160c0179efc902d4534e0}} \\
post-freeze target-only ablation & \texttt{\detokenize{238dbdc27ceb7ce016b9c995b5472cf1ad71a76632b6609b0df66408a1e26a7e}} \\
v91 legacy near-target-only confirmation & \texttt{\detokenize{3a0ee874afeff6df96a90fb4aed4e42f79705469f0d2a72da514932385990f67}} \\
v98 preselected controller-pair confirmation & \texttt{\detokenize{bb00da6d3ee0f9420bdd7827eeada7c0b4755fadae4a871e3de6c42ded1dc65d}} \\
v98 remote result identity audit & \texttt{\detokenize{f8243b8279e600721f978a763415ced656c3528b190fb952877b17f2b982a91b}} \\
Los Angeles deployment model & \texttt{\detokenize{f26ee8e5829e248ef411f77ad3e19ed0d63fd52cdfdae6d8291a9a3aeaaa7aff}} \\
Jinan deployment model & \texttt{\detokenize{a34970f1e4f023079650866097359ec3b2509e17787685249a21363cc7b03741}} \\
closed-loop source tree & \texttt{\detokenize{ad555ff72a11a62cc88c8bbbe2ab2f5df1db8de920af03e2f61d99aaba226a7d}} \\
\bottomrule
\end{tabular}
}
\end{table}

The publication artifacts are regenerated with
\begin{verbatim}
python3 paper_tsc/figures/build_tsc_external_confirmation.py
python3 paper_tsc/figures/build_tsc_causal_transfer_schematic.py
python3 paper_tsc/figures/build_tsc_postfreeze_stress_test.py
python3 paper_tsc/figures/build_tsc_source_gate_development.py
python3 paper_tsc/figures/build_tsc_source_remedy_ladder.py
\end{verbatim}
These commands form the figure and source-table submission contract.  The
unreferenced v16-era builder and its outputs are retained only as development
history and must not be substituted for Figures 1--3.
The first command refuses to write figures or tables if the development,
external offline, joint-freeze or closed-loop validity gate is not satisfied.
The third command separately requires all v86--v89 evidence-integrity audits
to pass and the preregistered v89 scientific decision to remain rejected.  The
fourth command requires the v145, v146 and v148 development decisions to remain
rejected and does not promote them to external confirmation.
The fifth command requires the canonical post-v148 remedy artifacts to retain
their recorded cohorts and adjudications, including the failed closed-loop
smoke, exact rigid fallbacks, the superseded V123/V157A comparison, the bounded
V157B B100 branch authorization and the incomplete V157C three-seed native
one-action diagnostic.  It also requires the V158 native-prefix OOF gate to
remain failed and its reserve authorization to remain false.
